%=====================================================================
\section{Questions and answers}
\label{sec:faq}

\newcommand{\qa}[2]{\par\noindent\textbf{Q: #1}\par\noindent #2\par\medskip}

\qa{What exactly is done in P4, and what is not?}{Every byte of every fingerprinted message is read by the P4
parser; the per-connection state, the decision what to copy, the JA4L times and the SSH counts are in the P4
pipeline. Outside P4: sorting, SHA-256, the SSH mode, writing the text, and joining split messages by framing
(Section~\ref{sec:glance}, D2--D6).}

\qa{Is the controller not doing the real work for split messages?}{No. It only puts bytes in order by their offsets
and cuts at framing boundaries (TLS record and handshake lengths, the empty line, the certificate list's lengths). It
then sends the whole message back, and the same P4 parser reads every value. A P4 switch cannot hold packets, so
some buffer must exist outside; the reading does not move with it.}

\qa{How does the switch handle ``the server's second packet arriving first''?}{From the SYN-ACK the switch knows the
server's initial sequence number, so it knows every segment's offset in the server's byte stream. A segment at a
non-zero offset in a direction that has not started is copied (rule (c), Fig.~\ref{fig:copyrules}); when the first
segment arrives, the controller has every piece and joins them (Section~\ref{sec:seq}). Tested in every order for 2
and 3 pieces, plus reversed 4 and 5 pieces.}

\qa{Is the server's side (h2) fingerprinted, or only the client's?}{Both. h1's ClientHello offers its versions,
ciphers, extensions and ALPN: JA4. h2 does not send a list of its own; its ServerHello states its choice (one
version, one cipher, the extensions it answers, the ALPN it picked): JA4S. In TLS 1.2 h2 then sends its certificate
chain in the clear: one JA4X per certificate (TLS 1.3 encrypts it). h2's SYN-ACK gives JA4TS, its half of the round
trip JA4L-S; JA4SSH counts both directions (\code{c\ldots} h1, \code{s\ldots} h2); h2's DHCP replies (Offer,
Advertise, Reply) get their own JA4D. Only HTTP has no server value: JA4+ fingerprints the request, not the
response. Live run, TLS 1.2: \code{t120300\_c030\_460f64128655} = TLS 1.2, 3 extensions, no ALPN, cipher
\code{c030} chosen from h1's 28.}

\qa{h2's messages over several segments, second half first?}{Handled the same way as the client's. From the
SYN-ACK the switch knows h2's initial sequence number, so every segment of h2 has a known offset. A ServerHello cut
in pieces is joined like a ClientHello (rules (a), (b), (c)). For the TLS 1.2 certificate flight: a later segment
that arrives before the ServerHello is copied by rule (c) (h2's direction has not started); the segment with the
ServerHello starts the flight (rule (d)); later segments are copied by rule (b); the controller cuts the
certificates out once there is no gap, and sends each back for the switch's DER walk. Tested:
\code{test\_server\_flight\_reorder}, 19 of 19 orders exact; JA4S also in the 202 split / reorder cases. Limit: a
ServerHello that is itself split starts no certificate copy, so JA4X is not reported (never wrong); a TLS 1.2
ServerHello is about 100 bytes, so TCP rarely splits it.}

\qa{Does the same hold for HTTP and SSH split over packets?}{HTTP: yes, for the client's requests (JA4+ has no
fingerprint of the server's response). A split request is joined like a hello: 76 of the 202 split / reorder cases
are HTTP. A \emph{later} request on a keep-alive connection, split and second piece first, needs rule (e): pieces
ahead of a gap in an HTTP direction are copied (18 of 18 orders exact), and the connection's JA4H are listed in
stream order. SSH: nothing is joined, by definition. JA4SSH counts TCP segments (payload lengths, packets, bare ACKs
per window of 200), as FoxIO's tool does; joining would make it differ. Counts do not depend on order, and the mode
takes the smallest length on a tie, so reordering inside a window changes nothing.}

\qa{What happens when a client reuses a port for a new connection?}{The SYN with a new sequence number gives the
connection a new number (ConnId); every control starts over and the controller opens a new record. An SSH
connection's last window is reported on that SYN. A repeated SYN (same sequence number) changes nothing. Before this
fix the second connection's JA4 was lost and two SSH sessions were added into one wrong JA4SSH
(\code{test\_port\_reuse}).}

\qa{Does it work beyond h1 -- s1 -- h2?}{The fingerprints do not depend on the topology, but some assumptions do:
both directions of a connection must pass the switch (asymmetric routing breaks JA4TS, JA4L, JA4S, JA4SSH and the
server-side join), one switch should fingerprint a connection, JA4L depends on where the switch is, and the program
is written for BMv2, not hardware. Port reuse, VLAN tags, SSH on other ports and bounded controller memory were
fixed after this review. Section~\ref{sec:assumptions} lists each assumption with its status.}

\qa{VLAN tags? SSH on another port?}{Up to two tags are read (802.1Q, and 802.1ad outside 802.1Q), as FoxIO's tool
reads them. SSH is recognised on any port by its ``SSH-'' banner, as tshark does under FoxIO's tool; bare ACKs
before the banner are counted, because FoxIO counts the handshake's ACK (\code{test\_vlan\_and\_ssh\_port}).}

\qa{And without the SYN?}{Then the offsets are learned from the first segment that starts a message; a piece that
came before it cannot be placed. The switch reports nothing rather than a wrong value (8 limit cases, tested).}

\qa{Why are packets resubmitted? Does that change them?}{The parser skips bytes with \code{advance}, which removes
them from the emitted packet. Resubmitting runs the original packet again; the second pass parses only Ethernet, so
all bytes are forwarded unchanged. Checked on every run, frame by frame (9,013 of 9,013).}

\qa{Why a hash of the endpoints and not a table entry per connection?}{The data plane cannot add table entries, and
waiting for the controller would lose the first packets (and JA4L's SYN time). Hashed register slots exist from the
first packet; a 32-bit check value detects another connection in the same slot (D8), and a connection number tells
a new connection on reused ports from the old one (D21).}

\qa{Why does the controller sort and hash instead of the switch?}{v1model has no SHA-256 and no loops in the
pipeline; the fingerprint is fixed by the values the switch read, so sorting and hashing them elsewhere gives the
same value. A sorting network in P4 would only be needed on hardware to emit sorted lists (D4).}

\qa{Why 128 ciphers, 64 extensions, 512 bytes of names \ldots?}{Chosen generously, then measured: all 1,238
lists of 42 captures fit (Section~\ref{sec:results}). Past a capacity the fingerprint is reported incomplete, never
truncated (tested).}

\qa{Why does a digest carry a serial number?}{BMv2 drops a digest identical to one not yet acknowledged; identical
DHCP retransmissions, repeated HTTP requests and identical SSH windows were lost. 32 bits is far more than the
number of unacknowledged digests can reach; wrapping is harmless because the value is never read (D10).}

\qa{Why is JA4L not scored exact?}{It is a time. BMv2 stamps a packet when its pipeline takes it up, which a
software switch delays (median 0.33\,ms in replay). The TTL part is scored; the latency is reported as a
distribution; the arithmetic is tested with known 1\,s / 0.5\,s gaps.}

\qa{How do you know the fingerprints are right?}{FoxIO's own tools compute the reference on the same frames:
573 exact, 5 differences explained by hand (the reference cannot decode those streams), 0 wrong, 0 missing, 0 extra;
the live Mininet run with real programs adds every fingerprint again: 49 of 49 exact, 196 of 196 frames forwarded
unchanged.}

\qa{What are the 5 EXPLAINED rows?}{One SOCKS4 ServerHello on port 9901 whose ClientHello is missing from the
capture, so FoxIO's tool decodes no TLS there (decoded by hand: equal to the switch). Four SSH streams that are
malformed or SSH version 1, where tshark refuses a few segments as SSH; the switch counts every port-22 payload
segment.}

\qa{Is fingerprinting complete in every respect?}{No, and the project does not claim it. What the evidence
supports: for clear-text TCP and DHCP traffic, every JA4+ fingerprint computable without keys matches FoxIO's tools,
for messages whole, split, reordered, duplicated or overlapping, from both sides (573 exact on FoxIO's samples,
47 of 47 in the final live run, 27 tests). Not done: JA4 over QUIC and DTLS, tunnels, IPv6 extension headers; two
cases are not tested (several Cookie headers, a HelloRetryRequest's second ClientHello). Section~\ref{sec:coverage}
gives each item with the reason it is not done and what doing it would take.}

\qa{What is not covered?}{Section~\ref{sec:limits}, known limits: encrypted content (HTTP/2, HTTP/3, TLS 1.3
certificates), QUIC and DTLS, tunnels, IPv6 extension headers, more than two VLAN tags; and the deployment
assumptions of Section~\ref{sec:assumptions}.}

\qa{Can a higher-level program drive the switch?}{Yes: through P4Runtime, as the controller does (write
\code{forward}, enable digests, send packet-out). The controller's Python class is the API; another P4Runtime client
can replace it.}

\qa{Where does each result file come from?}{\code{results/compare\_all.json}: \code{make compare};
\code{results/live/}: \code{make live} (which also writes \code{results\_live.tex}); \code{results/measure/}: \code{tools/measure\_field\_counts.py};
the PDFs' tables: \code{docs/results\_*.tex}, written from those files.}

\qa{How is this document kept current?}{Both PDFs include the same problems table (\code{docs/problems.tex}) and
results tables (\code{docs/results\_*.tex}). Every fix adds a row to the problems table; every design change updates
the matching section here. \code{make docs} builds both.}
